% Generated draft: requires mathematical and source review.
\documentclass[11pt]{article}
\usepackage[margin=25mm]{geometry}
\usepackage{fontspec}
\setmainfont{Latin Modern Roman}
\usepackage{amsmath,amssymb,mathtools,unicode-math}
\setmathfont{Latin Modern Math}
\usepackage{enumitem,needspace,xurl,hyperref}
\setlength{\parindent}{0pt}
\setlength{\parskip}{5pt}
\setlength{\emergencystretch}{3em}
\newcommand{\N}{\mathbb{N}}
\newcommand{\Z}{\mathbb{Z}}
\newcommand{\Q}{\mathbb{Q}}
\newcommand{\R}{\mathbb{R}}
\newcommand{\C}{\mathbb{C}}
\newcommand{\F}{\mathbb{F}}
\newcommand{\A}{\mathbb{A}}
\newcommand{\PP}{\mathbb{P}}
\newcommand{\E}{\mathbb{E}}
\newcommand{\eps}{\varepsilon}
\newcommand{\dd}{\,\mathrm{d}}
\begin{document}
\section*{UnsolvedMath ID: 30002983}
\textbf{Source problem code:} catalogue entry
\subsection{Definitions and mathematical statement}
Let \(S\) be a smooth projective surface, let \(\xi\) be a curve class on \(S\), and let \(g\) be the prescribed genus. The supplied record denotes by \(V_g^\xi\) a parameter space of reduced curves on \(S\) of class \(\xi\) and genus \(g\). It does not specify the base field, the equivalence relation defining the class, the construction of this parameter space, or the genus convention for reducible curves. In particular, reducedness must not silently be strengthened to irreducibility or connectedness.

Once those choices are fixed, put
\[
V_{g,\mathrm{nod}}^\xi
=\{[C]\in V_g^\xi:\text{every singularity of }C\text{ is an ordinary node}\}.
\]
For the geometric interpretation over a field \(k\), an ordinary node at a singular closed point \(p\) of \(C_{\bar k}\) means
\[
\widehat{\mathcal O}_{C_{\bar k},p}
\cong \bar k[[u,v]]/(uv),
\]
where \(\bar k\) is an algebraic closure of \(k\). Whether this geometric interpretation is intended requires confirmation. Under the convention that every singularity must be a node, smooth curves are included.

The density question, if Zariski density is intended, is whether
\[
\overline{V_{g,\mathrm{nod}}^\xi}^{\,V_g^\xi}=V_g^\xi.
\]
Equivalently, every nonempty relatively Zariski open subset of \(V_g^\xi\) should contain a curve whose singularities are all ordinary nodes. This formulation also covers the empty parameter space, for which the equality is vacuous. The record gives no further restrictions on \(S\), \(\xi\), or \(g\); whether the intended question is universal or restricted to particular families, including \(K3\) surfaces, remains unresolved.

The record also asks whether density can be proved using a proposed local criterion involving an equisingular ideal \(I\) and an object \(A\), expressed in the cleaned statement as
\[
I\subsetneq A
\quad\Longleftrightarrow\quad
\text{the singularity is worse than an ordinary node}.
\]
Neither \(A\), the precise equisingular ideal \(I\), nor the meaning of ``worse than an ordinary node'' is defined in the supplied record. Its original wording uses \(I\subset A\), whereas the cleaned wording explicitly requires proper containment. Accordingly, this criterion is retained as a source claim requiring clarification, not asserted as a theorem. The local-to-global deformation hypotheses needed to connect such a criterion to density are also unspecified.
\subsection{Short English statement}
The question asks whether reduced curves of a fixed class and genus on a smooth projective surface can be approximated within their parameter space by curves having only ordinary nodes. It also proposes an equisingular-ideal criterion as a possible route to density. The field, parameter space, genus convention, scope, and local ideal notation need clarification. The supplied registry labels the problem partially solved; this draft does not independently verify that assessment.
\subsection{Sources}
\begin{itemize}
\item Supplied local UnsolvedMath registry record, permanent ID \(30002983\), title \emph{Density of Nodal Curves in Equigeneric Families}, source problem code \texttt{catalogue entry}, in the Oberwolfach Reports Open Problems collection. No individual proposer is identified. The supplied proposed year is 2015.
\item Supplied report citation: \emph{Mini-Workshop: Singular Curves on \(K3\) Surfaces and Hyperkähler Manifolds} (2016), Oberwolfach Reports, pp.\ 2939--2967, DOI \texttt{10.4171/owr/2015/51}. \url{https://doi.org/10.4171/owr/2015/51}. The citation's 2016 date and the record's proposed year 2015 are preserved without resolving their relationship.
\item Supplied literature reference: Dedieu--Sernesi, \emph{Equigeneric and Equisingular Families}. \url{https://doi.org/10.5565/PUBLMAT_61117_07}.
\item Supplied literature reference: Dedieu--Sernesi, \emph{Density Problem}. \url{https://ems.press/content/serial-article-files/46597}.
\item Supplied literature reference: Wahl, \emph{Equisingular Deformations}. \url{https://doi.org/10.1090/S0002-9947-1974-0419439-2}.
\item Literature assessment dated exactly \texttt{2026-08-21}: the supplied curation labels the problem partially solved and reports equigeneric deformations to nodal curves in broad classes and some \(K3\) cases, together with counterexamples and remaining \(K3\) cases. It describes density as unsettled in general. This is supplied dated curation derived from the OWR record, not an independently verified literature assessment; the precise hypotheses and scope of those claims require review.
\item The registry assigns the statement status \texttt{corrected\_verified} and reports a check against the official Oberwolfach report. This is supplied source-verification metadata, not a verification performed for this draft. The background explicitly requests independent verification of time-sensitive status and cited claims.
\end{itemize}
\end{document}
